\section{Conclusion}

Completed development experiments distinguish nominal-family discrimination, practical predictive resolution and observation-design feasibility. A rich generator may be selected reliably while its predictive advantage remains smaller than the chosen computational tolerance; insufficient observations must remain visible in recovery denominators. These conclusions concern the specified synthetic benchmarks. They do not validate a human regulatory mechanism, and final inferential calibration remains separate from development evidence.

Mechanistic interpretation requires evidence that competing dynamical explanations can actually be distinguished by available observations. Preliminary ARC simulations show that P, PI, and PID generators can be difficult to recover even with known synthetic ground truth, and modest analyst measurement noise can make simpler baselines preferable. The revised study therefore treats observational equivalence as the primary scientific problem. Its goal is to identify when family recovery is reliable, when parameter estimates are practically informative, and when only predictive statements are defensible.
